\section{Results}
\subsection{
Static Translation using \dslr{}}
\label{sec:chamfer-nu}
CARLA-64 and ARD-16 have correspondence pairs available.
Hence, we report Chamfer's Distance and Earth Mover's Distance for them across the 5 baselines. (Table \ref{tab:carlaresults}). \dslr{}
performs better at both metrics against all the baselines.
% We also show ablation studies w.r.t. the components of our models over the best-performing  baseline in Table \ref{tab:ablation}. 
The key advantage of \dslr{} is that it is able to accurately differentiate between static and dynamic objects as well as between the corresponding occlusion patterns. It omits dynamic structures/occlusions and preserves static structures/occlusions (Video Demo at \cite{url1}). This effect is due to the explicit attention generated by the couple discriminator on dynamic objects and due to the strategy of contrasting LiDAR triplets using carefully chosen hard negatives. 
\setlength{\tabcolsep}{1.5pt}

 \begin{table*}[htbp]
\scriptsize
    \centering
    \begin{tabular}{|c|c|c|c|c|c|c|}

    \hline
     Run & C-Net         & S-INV        & C-AE         & C-GAN       &  Best-Baseline       & Ours        
\\ \hline
\multicolumn{7}{|c|}{\textbf{CARLA-64} (CD/EMD)} \\ \hline

 8 & 18.08/506        & 11.91/981  & 1.14/291.12  & 1.58/243.52 &  0.88/259.23 & \textbf{0.37/235.90} \\ \cline{1-7} 
 9 & 14.93/432     & 9.11/881   & 1.39/194.57  & 1.48/198.36 &  1.33/185.95 & \textbf{0.33/156.14}  \\ \cline{1-7} 
 10 & 17.43/692.32 & 10.23/887  & 6.06/483.59  & 5.11/463.12 &  1.57/231.37 & \textbf{1.19/226.81} \\ \cline{1-7} 
 11 & 6.20/608     & 8.56/981   & 1.92/324.26  & 1.59/292.21 &  1.24/239.97 & \textbf{0.95/207.79} \\ \cline{1-7} 
 12   & 17.91/597    & 9.32/901   & 3.89/356.20  & \textbf{3.54}/271.63 &  8.44/206.67 & 5.36/\textbf{186.35} \\ \cline{1-7} 
 13  &17.11/ 499.3  & 12.53/965  & 2.15/388.81  & 1.71/341.26 &  1.04/243.74 & \textbf{0.6/202.30}  \\ \cline{1-7} 
 14  &  16.18/501    & 14.12/1004 & 11.01/181.61 & 8.92/185.71 & \textbf{2.57}/228.34 & 2.63/\textbf{210.48}  \\ \hline
\multicolumn{7}{|c|}{\textbf{ARD-16} (CD/EMD)} \\ \hline

 1& 3.98/502     & 7.49/698  & 0.95/201  & 10.65/195 &  0.37/194 & \textbf{0.36/128} \\ \hline
\multicolumn{7}{|c|}{\textbf{KITTI} (LQI)} \\ \hline
% \MULTIROW{1}{*}{CARLA-64} 
 & 18.08/506        & 11.91/981  & 1.14/291.12  & 1.58/243.52 &  0.88/259.23 & \textbf{0.37/235.90} \\ \hline

    \end{tabular}
    \caption{Comparison of our model on static reconstructions with
5 baselines on 3 datasets. For all metrics lower is better. A video demo of the reconstructions is available at \cite{url1}.  }
    \label{tab:carlaresults}
    \end{table*}


We observe that our model works reasonably well even on
challenging sparse dataset - ARD-16. This is also highlighted by the fact that
\dslr{} is capable of generating precise dynamic object segmentation, unlike the best baseline (Figure \ref{fig:ati-moving-movable}) for ARD-16. \dslr{} overcomes the shortcoming of existing baselines by reconstructing only dynamic occlusion and retaining all other static structures. These static structures are vital for AVs in challenging scenarios like sparse datasets (ARD-16) on industrial floors.


For KITTI, we use LQI as a metric due to the unavailability of ground truth static scans. Our method improves to a large extent on the available baselines. The improvement in static translation translates to superior downstream  performance in segmentation-devoid settings (Table \ref{tab:slamresults}). 
% For a video demo of the static reconstruction results, click  \href{https://drive.google.com/file/d/1YDX4PBMQJMbuelSR_Ox7XjzXuOWu6bu_/view?usp=sharing}{here}.


\subsection{Segmentation of Dynamic Objects}
\label{sec:segm_mova_movi}

\begin{figure}[htbp]
  \centering
\includegraphics[width=0.95\linewidth, height=0.45\linewidth]{ARD-Movable.png}  
  \caption{Dynamic Object Segmentation comparison on ARD-16. \textbf{Left}: A corresponding LiDAR pair. Top is dynamic (\textcolor{red}{D}) and bottom is static (\textcolor{blue}{S}) scan. Note the presence of dynamic objects (\textcolor{red}{in red}) in the top dynamic scan (\textcolor{red}{D}). \textbf{Right}: Given dyanmic scan, \textcolor{red}{D} as input, segmentation output generated by our model is on top. Segmentation output by the best baseline is at bottom. \dslr{} precisely segments dynamic objects while the baseline is unable to capture them, which can be visualized in the Video demo \cite{url2}.} 

\label{fig:ati-moving-movable}
\end{figure}

\begin{figure}[t]
    \centering
    \includegraphics[width=6in]{moving-movable-carla.pdf}
    \caption{(Please view from Left to Right) \textbf{Row 1}: Top view of a contiguous  series of 4 dynamic scans at time intervals \textit{t, t+3, t+6, t+9} in a LiDAR sequence. Red denotes movable objects while blue denotes moving objects. 
    \textbf{Row 2}: When dynamic scan in row 1 is given as input to \dslr{}, it generates segmentation of movable objects (red) and moving objects (blue). Note that the moving object (\textcolor{blue}{in blue}) changes position across the scans from left to right while the movable objects (\textcolor{red}{in red}) are currently stationary. These movable objects, although stationary here, can be seen as moving in the video demo \cite{url3}.} 

    \label{fig:movable-second}
\end{figure}


We qualitatively show the segmentation of moving as well as movable objects in LiDAR scans (Figure \ref{fig:movable-first}, \ref{fig:ati-moving-movable} and \ref{fig:movable-second}) using \dslr{}. We achieve this by subtracting the original dynamic input from the generated static reconstruction (consisting of only static structures). The subtraction yields precise details of the non-stationary objects. These non-stationary objects can be
easily split into movable and movable by temporally observing the relative change in their position w.r.t the vehicle. The key advantage of \dslr{} is that unlike existing baselines, it accurately differentiates between dynamic and static obstacles and their patterns and removes only dynamic obstacles via static structure reconstruction. This results in accurate dynamic object segmentation (Figure \ref{fig:movable-second} and Video Demo at \cite{url3}). 
Moreover, challenging sparse datasets like ARD-16 have few points falling on dynamic objects. We show that the best baseline model fails to identify the dynamic objects for ARD-16. However, we are successful in segmenting out dynamic objects in such sparse datasets (Figure \ref{fig:ati-moving-movable} and Video Demo at \cite{url2}).





\subsection{Application of \dslr{} for navigation in dynamic environments}
In this section, we demonstrate the benefit of reconstructed static structures on navigation. Our model acts as a preprocessing step on LiDAR scans before they are consumed by the  navigation algorithm.

For KITTI, we use the KITTI Odometry dataset sequences from 00-10. For ARD-16, we use the single run released, which is longer and combines all sequences and scenes available in previous testing runs.  For CARLA-64, we use the 4 available dynamic LiDAR sequences.

We first compare \dslr{} against the best baselines in in the segmentation-free settings using all the three datasets in Table \ref{tab:slamresults}.

We also compare \dslruda{} with a segmentation-label assisted baseline - KITTI-Seg. KITTI-Seg utilizes semantic segmentation information to remove dynamic objects before using the LiDAR scans for navigation. We show these results on the KITTI dataset because of the availability of semantic segmentation information. The comparison of KITTI-Seg against \dslruda{} is available in Table \ref{tab:kitti-seg-all}.

\label{sec:slamindyenv}



\setlength{\tabcolsep}{1.5pt}
\begin{table}[t]
\scriptsize
\centering
\begin{tabular}{|ccccccc|}
\hline
\multicolumn{1}{|c|}{}           & \multicolumn{3}{c|}{\textbf{Best Baseline}}                                                                                            & \multicolumn{3}{c|}{\dslr{}}                                                                              \\ \hline
\multicolumn{1}{|c|}{Run}        & \multicolumn{1}{c|}{ATE}         & \multicolumn{2}{c|}{RPE}                                                          & \multicolumn{1}{c|}{ATE}            & \multicolumn{2}{c|}{RPE}                                          \\ \hline
\multicolumn{1}{|c|}{}           & \multicolumn{1}{c|}{(m)}         & \multicolumn{1}{c|}{Trans (m)}        & \multicolumn{1}{c|}{Rot ($\theta$)}         & \multicolumn{1}{c|}{(m)}            & \multicolumn{1}{c|}{Trans (m)}      & Rot ($\theta$)                \\ \hline
\multicolumn{1}{|c|}{}           & \multicolumn{6}{c|}{\textbf{CARLA-64}}                                                        \\ \hline
\multicolumn{1}{|c|}{0}          & \multicolumn{1}{c|}{5.67}        & \multicolumn{1}{c|}{0.102}            & \multicolumn{1}{c|}{0.413}                & \multicolumn{1}{c|}{\textbf{1.74}}  & \multicolumn{1}{c|}{\textbf{0.075}} & \textbf{0.412}              \\ \hline
\multicolumn{1}{|c|}{1}          & \multicolumn{1}{c|}{2.71}        & \multicolumn{1}{c|}{\textbf{0.050}}   & \multicolumn{1}{c|}{\textbf{0.390}}       & \multicolumn{1}{c|}{\textbf{1.10}}  & \multicolumn{1}{c|}{0.060}          & \textbf{0.390}              \\ \hline
\multicolumn{1}{|c|}{2}          & \multicolumn{1}{c|}{1.34}        & \multicolumn{1}{c|}{\textbf{0.077}}   & \multicolumn{1}{c|}{\textbf{0.526}}       & \multicolumn{1}{c|}{\textbf{1.11}}  & \multicolumn{1}{c|}{0.078}          & \textbf{0.520}              \\ \hline
\multicolumn{1}{|c|}{3}          & \multicolumn{1}{c|}{6.80}        & \multicolumn{1}{c|}{\textbf{0.136}}   & \multicolumn{1}{c|}{\textbf{0.375}}       & \multicolumn{1}{c|}{\textbf{3.83}}  & \multicolumn{1}{c|}{0.212}          & 0.379                       \\ \hline
\multicolumn{1}{|c|}{}           & \multicolumn{6}{c|}{\textbf{ARD-16}}                      \\ \hline
\multicolumn{1}{|c|}{0}          & \multicolumn{1}{c|}{2.06}        & \multicolumn{1}{c|}{\textbf{0.067}}   & \multicolumn{1}{c|}{1.550}                & \multicolumn{1}{c|}{\textbf{1.97}}  & \multicolumn{1}{c|}{\textbf{0.067}} & \textbf{1.520}                       \\ \hline
\multicolumn{1}{|c|}{}           & \multicolumn{6}{c|}{\textbf{KITTI}}                      \\ \hline

\multicolumn{1}{|c|}{0}          & \multicolumn{1}{c|}{5.22}        & \multicolumn{1}{c|}{0.025}            & \multicolumn{1}{c|}{0.005}                & \multicolumn{1}{c|}{\textbf{2.64}}  & \multicolumn{1}{c|}{\textbf{0.015}} & \textbf{0.003}              \\ \hline
\multicolumn{1}{|c|}{1}          & \multicolumn{1}{c|}{24.64}       & \multicolumn{1}{c|}{0.140}            & \multicolumn{1}{c|}{0.007}                & \multicolumn{1}{c|}{\textbf{4.91}}  & \multicolumn{1}{c|}{\textbf{0.080}} & \textbf{0.003}              \\ \hline
\multicolumn{1}{|c|}{2}          & \multicolumn{1}{c|}{8.98}        & \multicolumn{1}{c|}{0.032}            & \multicolumn{1}{c|}{0.003}                & \multicolumn{1}{c|}{\textbf{5.84}}  & \multicolumn{1}{c|}{\textbf{0.026}} & \textbf{0.002}              \\ \hline
\multicolumn{1}{|c|}{4}          & \multicolumn{1}{c|}{1.85}        & \multicolumn{1}{c|}{0.158}            & \multicolumn{1}{c|}{0.132}                & \multicolumn{1}{c|}{\textbf{0.77}}  & \multicolumn{1}{c|}{\textbf{0.108}} & \textbf{0.124}              \\ \hline
\multicolumn{1}{|c|}{5}          & \multicolumn{1}{c|}{2.05}        & \multicolumn{1}{c|}{0.026}            & \multicolumn{1}{c|}{\textbf{0.004}}       & \multicolumn{1}{c|}{\textbf{1.30}}  & \multicolumn{1}{c|}{\textbf{0.015}} & 0.005                       \\ \hline
\multicolumn{1}{|c|}{6}          & \multicolumn{1}{c|}{4.46}        & \multicolumn{1}{c|}{0.263}            & \multicolumn{1}{c|}{0.025}       & \multicolumn{1}{c|}{\textbf{1.82}}  & \multicolumn{1}{c|}{\textbf{0.060}} & \textbf{0.013}                       \\ \hline
\multicolumn{1}{|c|}{7}          & \multicolumn{1}{c|}{5.23}        & \multicolumn{1}{c|}{0.026}            & \multicolumn{1}{c|}{0.006}                & \multicolumn{1}{c|}{\textbf{2.64}}  & \multicolumn{1}{c|}{\textbf{0.016}} & \textbf{0.003}              \\ \hline
\multicolumn{1}{|c|}{8}          & \multicolumn{1}{c|}{6.70}        & \multicolumn{1}{c|}{0.031}            & \multicolumn{1}{c|}{\textbf{0.006}}       & \multicolumn{1}{c|}{\textbf{2.90}}  & \multicolumn{1}{c|}{\textbf{0.018}} & \textbf{0.006}              \\ \hline
\multicolumn{1}{|c|}{9}          & \multicolumn{1}{c|}{27.96}        & \multicolumn{1}{c|}{0.350}            & \multicolumn{1}{c|}{0.114}                & \multicolumn{1}{c|}{\textbf{4.50}}  & \multicolumn{1}{c|}{\textbf{0.055}} & \textbf{0.043}              \\ \hline
\multicolumn{1}{|c|}{10}         & \multicolumn{1}{c|}{2.45}        & \multicolumn{1}{c|}{0.034}            & \multicolumn{1}{c|}{0.010}                & \multicolumn{1}{c|}{\textbf{0.70}}  & \multicolumn{1}{c|}{\textbf{0.024}} & \textbf{0.008}              \\ \hline
\end{tabular}
\caption{Comparison of SLAM Results on 4 CARLA-64, 1-ARD and 10 KITTI Sequences without segmentation. For all metrics, lower is better.}
\label{tab:slamresults}
\end{table}




 \subsubsection{Analysis of navigation Results}

\textbf{Segmentation-devoid scenarios} - In Table \ref{tab:slamresults}, we demonstrate that in segmentation devoid settings, \dslr{} is able to improve navigation performance when compared against the best baseline over all the three datasets. This results deomonstrate the importance of the accurate static structures generated by \dslr{}. For KITTI we do not have corresponding static-dynamic LiDAR scans available. Inspite of this shortcoming, \dslruda{} is able to achieve superior results. This demonstrates that our approach can be successfully applied to datasets even in the absence of  correspondence information in label-devoid settings.


We notice that the improvement in navigation error for ARD-16 is marginal compared to the other datasets (Table \ref{tab:slamresults}). This is expected because ARD-16 is a sparse dataset with low density of points falling on the dynamic objects. The \textit{new static points} in the reconstructed static scan that replace the dynamic points are fewer in number.  The benefit on navigation is directly proportional
to the number of new static points introduced in the generated scan.
Therefore, the improvement on ARD-16 is expected to be less compared to the dense datasets. 
\begin{figure}[htbp]
% \centering
\begin{tabular}{|c|c|c|}
\hline
\subf{\includegraphics[width=0.33\linewidth]{00_trajectories_updated.png}}
     {Sequene 00}
&
\subf{\includegraphics[width=0.33\linewidth]{02_trajectories_updated.png}}
     {Sequene 02}

&
\subf{\includegraphics[width=0.33\linewidth]{08_trajectories_updated.png}}
     {Sequene 08}
\\

\hline
\subf{\includegraphics[width=0.33\linewidth]{01_trajectories_updated.png}}
     {Sequene 01}
&

\subf{\includegraphics[width=0.33\linewidth]{06_trajectories_updated.png}}
     {Sequene 06}
&
\subf{\includegraphics[width=0.33\linewidth]{09_trajectories_updated.png}}
     {Sequene 09}


\end{tabular}

\caption{Comparison of the estimated trajectory using \dslr{} with the best baseline for 6 sequences. \textbf{Row 1}- sequences with multiple loop closures. \textbf{Row 2} - Sequences with minimal (0 or 1)  loop closures.}
\label{fig:kitti-slam-trajectory}
\end{figure}

\begin{table}[]
\scriptsize
\centering

\begin{tabular}{|c|c|c|c|}
\hline
\multirow{2}{*}{Sequence} & \multirow{2}{*}{Loop Clousres} & \multirow{2}{*}{\begin{tabular}[c]{@{}c@{}}Dynamism Level\\ (Number of dynamic objects)\end{tabular}} & \multirow{2}{*}{\begin{tabular}[c]{@{}c@{}}Lenght\\ (Number of LiDAR scans)\end{tabular}}  \\
                          &                                &             &                                                                                          \\ \hline
0                         & \textgreater{}=2               & Medium(21)  &        4390                                                                             \\ \hline
1                         & 1                              & High(53)    &        1173                                                                                \\ \hline
2                         & \textgreater{}=2               & Low(9)      &        4485                                                                               \\ \hline
4                         & 0                              & High(30)    &        284                                                                              \\ \hline
5                         & \textbf{\textgreater{}=2}               & \textbf{High(40)}    &        \textbf{2750}                                                                             \\ \hline
6                         & 1                              & Low(10)     &        1091                                                                            \\ \hline
7                         & 1                              & Medium(18)  &        1111                                                                           \\ \hline
8                         & \textbf{\textgreater{}=2}               & \textbf{High(49)}    &        \textbf{5149}                                                                          \\ \hline
9                         & 1                              & Low(13)     &        1599                                                                         \\ \hline
10                        & 0                              & Low(10)     &        1215                                                                        \\ \hline
\end{tabular}
\caption{Variation observed in the KITTI sequences w.r.t. number of loop closures and level of dynamism - associated with number of dynamic objects and length of sequence.}
\label{tab:kitti-variation}
\end{table}



We compare the estimated trajectories of \dslr{} with the best baseline on the urban KITTI Odometry sequence.
The KITTI Odometry sequences are diverse on several parameters - loop closures, length, dynamic crowded urban sequences to sparsely dynamic highway sequences. We show the variations in Table \ref{tab:kitti-variation}. Sequences 0, 2, 5, 8 have multiple loops closure while the other sequences have less than 2 loop closures. Sequences 1, 4, 5, 8 have high dynamism compared to the other sequences.



We compare the estimated trajectories of \dslr{} with the best baseline in Figure \ref{fig:kitti-slam-trajectory}. We observe a trend across the estimated  trajectories for the baseline - the navigation accumulates translation errors consistently as the sequence progresses despite loop closures. The top row in Figure \ref{fig:kitti-slam-trajectory} compares sequences with multiple loop closures. The baseline (indicated by green line) continues to accumulate translation error despite multiple loop closures as the run progresses from the start point (indicated by black dot). \dslr{} (indicated by red) is able to narrow down the error on these sequences. We observe a similar behaviour for sequences with minimal to no loop closures (row 2 of Figure \ref{fig:kitti-slam-trajectory}).



We investigate the reason behind the performance gain of \dslr{} using the KITTI dataset in Figure \ref{fig:kitti-slam-xyz-analysis}. We  plot the translation pose estimates of the best baseline and \dslr{} \textit{w.r.t.} the ground truth pose estimates for the KITTI sequences. Our analysis suggests that the baseline faces difficulty in accurately regressing the z-values for LiDAR scans during static reconstructions while leads to high pose error along the z-axis. Notably, it has significant errors along the z-axis, in addition to errors in the x and y axes. \dslr{} reduces the pose error along the z-axis and also reduces the error along the x and y axes for most sequences.


\begin{figure}

% \centering
\begin{tabular}{|c|c|c|}
\hline
\subf{\includegraphics[width=0.33\linewidth]{00_xyz_view.png}}
     {Sequene 00}
&
\subf{\includegraphics[width=0.33\linewidth]{01_xyz_view.png}}
     {Sequene 01}

&
\subf{\includegraphics[width=0.33\linewidth]{02_xyz_view.png}}
     {Sequene 02}
\\
\hline
\subf{\includegraphics[width=0.33\linewidth]{04_xyz_view.png}}
     {Sequene 04}
&
\subf{\includegraphics[width=0.33\linewidth]{05_xyz_view.png}}
     {Sequene 05}

&
\subf{\includegraphics[width=0.33\linewidth]{06_xyz_view.png}}
     {Sequene 06}
\\

\hline
\subf{\includegraphics[width=0.33\linewidth]{07_xyz_view.png}}
     {Sequene 07}
&
\subf{\includegraphics[width=0.33\linewidth]{09_xyz_view.png}}
     {Sequene 09}

&
\subf{\includegraphics[width=0.33\linewidth]{10_xyz_view.png}}
     {Sequene 10}
\\
\hline


\end{tabular}

\caption{ SLAM error along the x, y, z coordinates for KITTI sequences. Blue denotes GT, Red denotes Ours while Green Denotes Best baseline. We observe that the best baseline constantly has high errors along the z-dimension as compared to MOVES for most of the sequences. Refer to Table 2 for a quantitative analysis of the translation component of navigation error.}
\label{fig:kitti-slam-xyz-analysis}
\end{figure}

\begin{figure}

% \centering
\begin{tabular}{|c|c|c|}
\hline
\subf{\includegraphics[width=0.33\linewidth]{00_rpy_view.png}}
     {Sequene 00}
&
\subf{\includegraphics[width=0.33\linewidth]{01_rpy_view.png}}
     {Sequene 01}

&
\subf{\includegraphics[width=0.33\linewidth]{02_rpy_view.png}}
     {Sequene 02}
\\
\hline
\subf{\includegraphics[width=0.33\linewidth]{04_rpy_view.png}}
     {Sequene 04}
&
\subf{\includegraphics[width=0.33\linewidth]{05_rpy_view.png}}
     {Sequene 05}

&
\subf{\includegraphics[width=0.33\linewidth]{06_rpy_view.png}}
     {Sequene 06}
\\

\hline
\subf{\includegraphics[width=0.33\linewidth]{07_rpy_view.png}}
     {Sequene 07}
&
\subf{\includegraphics[width=0.33\linewidth]{09_rpy_view.png}}
     {Sequene 09}

&
\subf{\includegraphics[width=0.33\linewidth]{10_rpy_view.png}}
     {Sequene 10}
\\
\hline


\end{tabular}

\caption{ Navigation error along the rotational pose estimates (roll, pitch, yaw) for KITTI sequences. Blue denotes GT, Red denotes Ours while Green Denotes Best baseline. \dslr{} estimates the rotational poses closer to the ground truth consistently across the KITTI sequences. Refer to Table 2 for a quantitative analysis of rotational component of navigation error.}
\label{fig:kitti-slam-rpy-analysis}
\end{figure}



We also investigate the rotational pose estimates (roll, pitch, yaw) for the KITTI sequences and compare the performance of the baseline against \dslr{} in Figure \ref{fig:kitti-slam-rpy-analysis}. We observe that although the baseline works on par with \dslr{} for the yaw estimates, it performs sub-optimally for roll and pitch estimates across the KITTI sequences. \dslr{} consistently has lower Relative Pose error compared to the baseline (Table \ref{tab:slamresults} and Figure \ref{fig:kitti-slam-rpy-analysis}).


\textbf{Comparison of segmentation-assisted method (KITTI-Seg) against MOVES-MMD} - We compare the navigation performance of \dslruda{} (which works without segmentation assistance), against KITTI-Seg that utilizes segmentation information to remove dynamic objects, in Table \ref{tab:kitti-seg-all}. 

 \dslruda{} works without segmentation labels. As expected, it is not always superior compared to KITTI-Seg.
We surprisingly observer that for sequence 5, 8, \dslruda{} performs superior to KITTI-Seg on ATE without the availability of segmentation labels. These sequences have very \textit{high dynamism}, are relatively \textit{longer} and hence have \textit{multiple loop closures} unlike other sequences (Table \ref{tab:kitti-variation}). These are very practical situation that most AVs come across frequently in everyday scenarios. Our model excels in these settings. Our investigation suggest that sequences with high dynamism have fewer static structures. KITTI-Seg removes dynamic objects and occlusions. The LiDAR scans are left with the few static structures. In such cases, KITTI-Seg pose estimates are not robust to errors accumulated along the sequence, especially at turns where even consecutive static scans have relatively less overlap.  On the contrary, \dslruda{} not only removes dynamic objects and occlusion but also introduces multiple new static structures that replace the dynamic objects and corresponding occlusions. This leads to an overall increase in the amount of static content in LiDAR scans. These additions lead to robust pose estimation as opposed to \textit{highly dynamic settings} that have fewer static structures. Thus, the LiDAR estimated poses are highly resistant to drift based errors (that accumulate mostly at turns) due to robust scan matching using multiple static structures. This robustness is reinforced by the presence of several loop closures that ensure that the accumulated drift is minimized. Thus, longer KITTI sequences with high dynamism and multiple loop closures demonstrate better navigation performance using \dslruda{}, even in the absence of segmentation information.
\setlength{\tabcolsep}{0.8pt}
\begin{table}[H]
\scriptsize
\centering
\begin{tabular}{|c|c|c|c|c|}
\hline
\textbf{Run}                 & \textbf{Model}                        & \textbf{ATE}            & \multicolumn{2}{|c|}{{\textbf{RPE}}}   \\ \hline
                 &                         &       \textbf{Trans(m)}      &  \textbf{Trans(m)}      & \textbf{Rot($\theta$)}        \\ \hline
\multirow{2}{*}{0}  
& KITTI-Seg            &     \textbf{0.33}          & \textbf{0.008}          & 0.013          \\  
& \dslruda{} (No Seg)              & 2.64  & 0.015 & \textbf{0.003}  \\ 
\hline
\multirow{2}{*}{1}  
& KITTI-Seg            & \textbf{0.98}         & 0.024          & \textbf{0.002}          \\  
& \dslruda{} (No Seg)              &     4.91   &  \textbf{0.08}   &    0.003 \\\hline

\multirow{2}{*}{2} 
& KITTI-Seg            & \textbf{3.74}         & \textbf{0.008}          & 0.004           \\  
& \dslruda{} (No Seg)              & 5.84   & 0.026 & \textbf{0.002} \\\hline

\multirow{2}{*}{4}
& KITTI-Seg            &         \textbf{0.38} & \textbf{0.085}          & \textbf{0.070}          \\  
& \dslruda{} (No Seg)              & 0.77           & 0.108 & 0.124 \\\hline
\multirow{2}{*}{5}

& KITTI-Seg            &     1.85          & \textbf{0.005}          & \textbf{0.003}          \\  
& \dslruda{} (No Seg)              & \textbf{1.30} & 0.015 & 0.005 \\\hline
\multirow{2}{*}{6}  
& KITTI-Seg            &        \textbf{0.30}         &    \textbf{0.012}        &   \textbf{0.002}        \\  
& \dslruda{} (No Seg)              & 1.82           & 0.06  & 0.013 \\\hline
\multirow{2}{*}{7}  

& KITTI-Seg            &       \textbf{0.42}          & \textbf{0.005}          & 0.004          \\  
& \dslruda{} (No Seg)              & 2.64           & 0.016 & \textbf{0.003} \\\hline

\multirow{2}{*}{8}  

& KITTI-Seg            &       81.82          &   0.316          & 0.135          \\  
& \dslruda{} (No Seg)              & \textbf{2.64}           & \textbf{0.016} &         \textbf{0.003} \\\hline

\multirow{2}{*}{9} 
& KITTI-Seg            &       \textbf{1.35}         & \textbf{0.005}          & \textbf{0.003}          \\  
& \dslruda{} (No Seg)              & 4.50       & 0.055 &    0.043 \\\hline
\multirow{2}{*}{10}
& KITTI-Seg            &     \textbf{0.39}          & \textbf{0.009}           & \textbf{0.006}          \\  
& \dslruda{} (No Seg)              & 0.70   & 0.024 & 0.008 \\\hline
\end{tabular}
\caption{Detailed Comparison of MOVES (without segmentation) against segmentation baselines (KITTI-Seg) on KITTI sequences.}
\label{tab:kitti-seg-all}
\end{table}




\subsection{Ablation Studies}
\label{sec:ablation}

 \begin{table}[h]
 \scriptsize
    \centering
    \begin{tabular}{|c|c|c|c|}
    \hline
    Run & Best Baseline & CoD & CoD+CL \\ \hline
    & \multicolumn{3}{|c|}{{CD/EMD}}      \\ \hline

    8   & 0.88/259 & 0.52/\textbf{232}   & \textbf{0.38}/236  \\ \hline
    9   & 1.33/186  & 0.50/164  & \textbf{0.30}/\textbf{156}\\\hline

    10  & 1.57/231 & 1.26/\textbf{221}   &  \textbf{1.16}/227 \\ \hline
    11  & 1.24/240 & 1.04/220   & \textbf{0.94}/\textbf{208} \\ \hline
    12  & 8.44/207 & 5.08/195   & \textbf{5.07}/\textbf{186}  \\ \hline
    13  & 1.04/244 & 0.65/204   & \textbf{0.60}/\textbf{202}  \\ \hline 
    14  & 2.57/228 & \textbf{2.50}/218 & \textbf{2.50}/\textbf{210}\\ \hline 
    % 15  & 1.37& 0.52& \textbf{0.31}  \\ \hline
    \end{tabular}
    \caption{Detailed ablation studies to show the effect of the couple discriminator (CoD) and contrastive loss (CL) on \dslr{} when compared to best baseline on CARLA-64 test set.}
    \label{tab:ablation}
    \end{table}





\setlength{\tabcolsep}{1.5pt}
\begin{table}[htbp]
\scriptsize
\centering
\begin{tabular}{|c|cc|}
\hline
Dataset                   & \multicolumn{1}{c|}{\begin{tabular}[c]{@{}c@{}}Vanilla Discriminator based GAN\end{tabular}} & \dslr{}                \\ \hline
                          & \multicolumn{2}{c|}{CD/EMD}                                                                                           \\ \hline
\multirow{7}{*}{CARLA-64} & \multicolumn{1}{c|}{1.58/243.52}                                                               & \textbf{0.37/235.90} \\ \cline{2-3} 
                          & \multicolumn{1}{c|}{1.48/198.36}                                                               & \textbf{0.33/156.14} \\ \cline{2-3} 
                          & \multicolumn{1}{c|}{5.11/463.12}                                                               & \textbf{1.19/226.81} \\ \cline{2-3} 
                          & \multicolumn{1}{c|}{1.59/292.21}                                                               & \textbf{0.95/207.79} \\ \cline{2-3} 
                          & \multicolumn{1}{c|}{\textbf{3.54}/271.63}                                                               & 5.36\textbf{/186.35} \\ \cline{2-3} 
                          & \multicolumn{1}{c|}{1.71/341.26}                                                               & \textbf{0.6/202.30}  \\ \cline{2-3} 
                          & \multicolumn{1}{c|}{8.92/\textbf{185.71}}                                                               & \textbf{2.63/}210.48 \\ \hline
\end{tabular}



\caption{Ablation study to show the comparison between a vanilla discriminator based GAN and our couple discriminator based \dslr{}.}
\label{tab:ablation-disc}
\end{table}

\dslr{} has 2 modules: Couple Discriminator(CoD) and Contrastive Loss(CL). We perform ablation studies to show the effect of these modules on \dslr{}. CoD is our couple-discriminator based Relativistic Average Least Squares GAN (Column 3 of Table \ref{tab:ablation}). It bridges a significant gap between the best baseline and our model. CL coupled with CoD (Column 4) gives further improvement on Chamfer's Distance. 

We also conduct another ablation study to demonstrate that a standard vanilla discriminator based Relativistic Average Least Square GAN performs inferior to our couple discriminator based \dslr{}. (Table \ref{tab:ablation-disc}). 


